For subgroups \(P_i\leq G\), typed coset Helly means that every pairwise
intersecting family of right cosets \(P_i g_i\), with distinct indices,
has a common element. This property is separate from IP. It supplies a
common representative for an incident flag, whereas IP determines the
stabilizer of a flag already represented in this way.

\textbf{Lemma UD.Helly (anchored criterion).} A finite indexed subgroup
family \((P_i)_{i\in I}\) has typed coset Helly if and only if the
following condition holds. For every \(i\in I\), every
\(J\subseteq I\setminus\{i\}\) with \(|J|\geq2\), and any selected
\(j_0\in J\), each \(g\in P_{j_0}\) satisfying \(g\in P_iP_j\) for all
\(j\in J\) admits \(k\in\bigcap_{j\in J}P_j\) with \(P_i g=P_i k\).

\emph{Proof.} The cosets \(P_i g\) and \(P_j\) intersect precisely when
\(g\in P_iP_j\). Thus Helly gives a common element \(k\), proving
necessity. Conversely, induct on the number of cosets in a pairwise
intersecting family. A family of at most two already has a common
element. For a larger family, the induction hypothesis gives a common
element of all but one coset. Right multiplication by its inverse turns
those cosets into the base cosets \(P_j\). The remaining coset meets
\(P_{j_0}\), so it has a representative \(g\in P_{j_0}\). Its
intersections with the other base cosets are precisely the product
conditions in the statement. The criterion supplies \(k\), and undoing
the right translation gives a common element of the original family.
\(\square\)

\section{Panel systems with a nontrivial chamber
stabilizer}\label{panel-systems-with-a-nontrivial-chamber-stabilizer}

\textbf{Lemma UD.CosetRealization.} Let a group \(G\) act
type-preservingly and flag-transitively on an incidence geometry of
finite rank, and let \((x_i)_{i\in I}\) be a base chamber. Write
\(H_i=\operatorname{Stab}_G(x_i)\) and use a right action. Then \[
H_i g\longmapsto x_i^g
\] is a type-preserving \(G\)-equivariant isomorphism from the subgroup
coset system to the given geometry.

\textbf{Proof.} Transitivity on one-object flags makes the map onto each
type. It is well defined and injective because \(x_i^g=x_i^h\) holds
exactly when \(gh^{-1}\in H_i\), which is equivalent to \(H_i g=H_i h\).
If two cosets share a representative \(k\), their images are
\(x_i^k,x_j^k\) and are incident, as the base objects are incident.
Conversely, if two images are incident, transitivity on flags of their
two types supplies a \(k\) carrying the base pair to them. The injective
object correspondence then writes both cosets with that common
representative. The map therefore preserves and reflects incidence.
\(\square\)

Let \(G\) be finite, let \(I\) be a finite type set, and let each panel
subgroup \(P_i\) contain \(B\). Put \(K_\varnothing=B\) and
\(K_S=\langle P_i:i\in S\rangle\) for nonempty \(S\), with \(K_I=G\).
This marking by subgroups is different from a marking by individual
involutions.

\textbf{Theorem UD.PanelRecognition.} Suppose \[
K_S\cap K_T=K_{S\cap T}\quad(S,T\subseteq I),\qquad [P_i:B]\ge2,
\tag{E4}
\] and suppose the maximal parabolics \(H_i=K_{I\setminus\{i\}}\) have
typed coset Helly. Then their coset incidence system is a firm
flag-transitive residually connected geometry with chamber stabilizer
\(B\). Conversely every finite-rank firm flag-transitive residually
connected coset geometry is obtained this way by taking the panel
stabilizers.

\textbf{Proof.} Equation (E4) identifies the stabilizer of every partial
base flag with the corresponding \(K_S\). UD.Helly supplies a common
group element in the cosets of every flag, and hence flag transitivity
and identification of the residues with the corresponding coset systems
in \(K_S\). Each rank-one residue has \([P_i:B]\) elements, giving
firmness. A rank-two residue is connected because its two panel
subgroups generate \(K_{\{i,j\}}\). For larger rank, all panel subgroups
occur in proper deletions, so the residue group is generated by its
maximal parabolics. The usual coset-path argument proves connectedness:
paths multiply successive elements of those subgroups, and conversely
their products give paths.

Conversely let \(K_S\) stabilize the complementary base flag. Stabilizer
intersections give (E4). In a connected residue, the subgroup generated
by its maximal parabolics has exactly the orbit reachable by coset
paths, so it is the whole residue group. Induction on \(|S|\)
reconstructs \(K_S\) from the panels, starting with \(K_\varnothing=B\).
Flag transitivity gives typed Helly by moving a flag to the base flag.
\(\square\)

The source's remaining axioms are separate, exact local tests. RWPRI
means that for every nonempty \(S\) there is an \(i\in S\) with
\(K_{S\setminus\{i\}}\) maximal proper in \(K_S\). Local
two-transitivity means that \(P_i\) has two orbits on ordered pairs in
\(B\backslash P_i\); equivalently \(B\) has two orbits on that coset
set. The source's geometric \((\mathrm{IP})_2\) means that every
rank-two residue is a partial linear space or a generalized digon. It is
not the subgroup identity (E4).

For the last test let \(P,Q\) be two panels, \(P\cap Q=B\) and
\(K=\langle P,Q\rangle\). The coset graph is a digon exactly when
\(|P||Q|=|K||B|\). Otherwise local two-transitivity makes a single
common-neighbour count sufficient. Choose \(x\in Q\setminus B\).
Neighbours of \(P\) are the \(Q\)-cosets represented by a right
transversal of \(B\) in \(P\); neighbours of \(Px\) are their translates
by \(x\). The stabilizer \(Q\) is two-transitive on its incident
\(P\)-type points, so all distinct points having a common \(Q\)-type
neighbour have the same number of common neighbours. This number must be
one. Pairs with no common neighbour already satisfy partial linearity.
This argument also applies when a local action has a nontrivial kernel.

\subsection{Complete augmentation through maximal
residues}\label{complete-augmentation-through-maximal-residues}

\textbf{Lemma UD.Transport.} Suppose representatives of all ambient
\(G\)-conjugacy classes of subgroups \(K\) and all \(K\)-conjugacy
classes of maximal subgroups \(H_0<K\) are given. For a specified \(H\),
choose any transporter \(c\) with \(H_0^c=H\) when one exists. All
maximal overgroups of \(H\) in that inclusion class are precisely
\((K^c)^n\) with \(n\in N_G(H)\).

\textbf{Proof.} Conjugate any maximal inclusion \(H<L\) to a
representative \(K\) of the class of \(L\). Its image of \(H\) is
\(K\)-conjugate to one of the listed \(H_0\). Comparing the resulting
transporter with \(c\) leaves an element normalizing \(H\). Conversely
each displayed conjugate contains \(H\) maximally. Taking the union over
all inclusion classes gives every maximal overgroup. \(\square\)

Let \(\mathcal P(B)\) be the finite set of overgroups \(P\) of \(B\)
such that \(B\) is maximal proper in \(P\) and \(P\) acts
two-transitively on \(B\backslash P\). UD.Transport computes this entire
set by retaining only the relevant local inclusions. The normalizer
\(N_G(B)\) acts on \(\mathcal P(B)\), and its permutation image has
exactly the same configuration orbits. We retain the actual subgroups
behind the permutation labels, so kernel reduction does not change any
geometric test.

\textbf{Theorem UD.MaximalAugmentation.} Start with all rank-one panels
modulo \(N_G(B)\). For every accepted panel tuple, let \(H\) be its
generated join. For every maximal overgroup \(K\) of \(H\), try every
\(P\in\mathcal P(B)\) with \(P\le K\) and \(P\nsubseteq H\), modulo the
pointwise stabilizer of the old panel tuple. Apply UD.PanelRecognition
and all the local source tests, then remove duplicates under \(N_G(B)\)
and type permutations. Iterating until the next level is empty gives
every target geometry with Borel \(B\) and acting group a subgroup of
\(G\).

\textbf{Proof.} For a target rank-\(r\) geometry, RWPRI supplies at
least one maximal parabolic \(H\) in its acting group \(K\). Deleting
its complementary panel gives a rank-\((r-1)\) residue with the same
Borel and all inherited source axioms. Induction puts its panel tuple in
a previously enumerated orbit. Its remaining panel lies inside \(K\) and
outside \(H\), and UD.Transport includes \(K\) among the maximal
overgroups to be tried. UD.ExtensionOrbits shows that taking
pointwise-stabilizer representatives loses no extension; a smaller
stabilizer than the full setwise one merely permits duplicates. Final
orbit tests remove those duplicates. Conversely every accepted tuple
satisfies exactly the recognition and local conditions.

By (E4) and proper panels, adjoining a new type strictly enlarges the
generated join: if \(P_i\) lay in \(K_S\) for \(i\) outside \(S\), then
\(P_i\cap K_S=B\) would force \(P_i=B\). Hence every tuple gives a chain
of strict subgroup inclusions from \(B\). Its rank is at most
\(\Omega([G:B])\), the number of prime factors with multiplicity. Thus a
run through the next level must reach an empty level. The certificate
requires that emptiness, not merely a chosen numeric cap. Finally take
\(B\) over all \(G\)-conjugacy classes of subgroups. \(\square\)
